\newpage
\section*{Additional assigned exercises}
\noindent These entries complete the remaining problems specified on the syllabus. Graph-matching answers identify each surface by its defining geometric features.

\subsection*{12.3 The Dot Product}
\begin{enumerate}[label=\textbf{\arabic*.}, start=27]
\item $(\ihat+\jhat)\times(\ihat+\khat)=\langle1,-1,-1\rangle$, so one suitable unit vector is \[\answer{\frac{\langle1,-1,-1\rangle}{\sqrt3}}.\]
\setcounter{enumi}{32}
\item For $\vect v=\langle4,1,8\rangle$, $|\vect v|=9$. Thus the direction cosines are \[\answer{\left(\frac49,\frac19,\frac89\right)},\] and the direction angles are \[\answer{\alpha\approx63.6^\circ,\quad\beta\approx83.6^\circ,\quad\gamma\approx27.3^\circ}.\]
\setcounter{enumi}{40}
\item With $\vect a=\langle4,7,-4\rangle$ and $\vect b=\langle3,-1,1\rangle$, $\vect a\cdot\vect b=1$ and $|\vect a|=9$. Hence \[\answer{\operatorname{comp}_{\vect a}\vect b=\frac19,\qquad \operatorname{proj}_{\vect a}\vect b=\frac{\langle4,7,-4\rangle}{81}}.\]
\setcounter{enumi}{46}
\item $\operatorname{comp}_{\langle3,0,-1\rangle}\vect b=2$ requires $\langle3,0,-1\rangle\cdot\vect b=2\sqrt{10}$. For example, \[\answer{\vect b=\left\langle\frac{2\sqrt{10}}3,0,0\right\rangle}.\]
\setcounter{enumi}{52}
\item Let $\vect n=\langle a,b\rangle$, a normal to $ax+by+c=0$. Projecting the vector from any point on the line to $P_1$ onto $\vect n$ gives \[d=\frac{|ax_1+by_1+c|}{\sqrt{a^2+b^2}}.\] For $(-2,3)$ and $3x-4y+5=0$, \[\answer{d=\frac{|-6-12+5|}{5}=\frac{13}{5}}.\]
\end{enumerate}

\subsection*{12.4 The Cross Product}
\begin{enumerate}[label=\textbf{\arabic*.}, start=1]
\item $\langle2,3,0\rangle\times\langle1,0,5\rangle=\answer{\langle15,-10,-3\rangle}$; its dot product with each original vector is $0$.
\setcounter{enumi}{4}
\item \[\left\langle\frac12,\frac13,\frac14\right\rangle\times\langle1,2,-3\rangle=\answer{\left\langle-\frac32,\frac74,\frac23\right\rangle}.\]
\setcounter{enumi}{6}
\item \[\langle t^3,t^2,t\rangle\times\langle t,2t,3t\rangle=\answer{\langle3t^3-2t^2,\ t^2-3t^4,\ 2t^4-t^3\rangle}.\]
\setcounter{enumi}{8}
\item $(\ihat\times\jhat)\times\khat=\khat\times\khat=\answer{\vect0}$.
\setcounter{enumi}{10}
\item $(\jhat-\khat)\times(\khat-\ihat)=\answer{\ihat-\jhat+\khat}$.
\setcounter{enumi}{14}
\item The diagram gives $|\vect u|=4$, $|\vect v|=3$, and $\theta=150^\circ$. Therefore $|\vect u\times\vect v|=12\sin150^\circ=\answer6$; the right-hand rule gives \answer{\text{into the page}}.
\setcounter{enumi}{16}
\item \[\vect a\times\vect b=\langle2,-1,3\rangle\times\langle4,2,1\rangle=\answer{\langle-7,10,8\rangle},\qquad \vect b\times\vect a=\answer{\langle7,-10,-8\rangle}.\]
\setcounter{enumi}{18}
\item \[\langle3,2,1\rangle\times\langle-1,1,0\rangle=\langle-1,-1,5\rangle,\] so the two unit vectors are \[\answer{\pm\frac{\langle-1,-1,5\rangle}{3\sqrt3}}.\]
\setcounter{enumi}{28}
\item $\overrightarrow{PQ}=\langle2,1,3\rangle$ and $\overrightarrow{PR}=\langle5,4,2\rangle$, so a normal is $\langle-10,11,3\rangle$ and \[\answer{\operatorname{Area}(\triangle PQR)=\frac{\sqrt{230}}2}.\]
\setcounter{enumi}{32}
\item $|\vect a\cdot(\vect b\times\vect c)|=|\langle1,2,3\rangle\cdot\langle2,8,-3\rangle|=\answer9$.
\setcounter{enumi}{36}
\item $\vect v\times\vect w=\langle4,12,32\rangle$, and $\vect u\cdot(\vect v\times\vect w)=4+60-64=\answer0$; hence the vectors are coplanar.
\setcounter{enumi}{42}
\item Since $|\vect a\times\vect b|=3$ and $\vect a\cdot\vect b=\sqrt3$, $\tan\theta=3/\sqrt3=\sqrt3$, giving \answer{\theta=60^\circ}.
\setcounter{enumi}{43}
\item Solving $\langle1,2,1\rangle\times\langle p,q,r\rangle=\langle3,1,-5\rangle$ gives \[\answer{\vect v=\langle1,-3,0\rangle+t\langle1,2,1\rangle}.\] No vector can give $\langle3,1,5\rangle$, because that vector has dot product $10$ with $\langle1,2,1\rangle$, whereas every cross product with it must be orthogonal to it.
\end{enumerate}

\subsection*{12.5 Equations of Lines and Planes}
\begin{enumerate}[label=\textbf{\arabic*.}, start=3]
\item \[\answer{\vect r=\langle-1,8,7\rangle+t\langle6,4,3\rangle};\quad x=-1+6t,\ y=8+4t,\ z=7+3t.\]
\setcounter{enumi}{4}
\item A normal is $\langle3,-2,2\rangle$, so \[\answer{3x-2y+2z=3}.\]
\setcounter{enumi}{14}
\item \[\answer{\frac{x-1}{-1}=\frac{y+5}{2}=\frac{z-6}{-3}}.\] Its intercepts are $(0,-3,3)$, $(-3/2,0,-3/2)$, and $(-1,-1,0)$.
\setcounter{enumi}{18}
\item Comparing the first two coordinate equations yields incompatible values for $t-2s$; therefore the lines are \answer{\text{skew}}.
\setcounter{enumi}{20}
\item Writing the symmetric equations parametrically and solving gives $t=2$, $s=1$, hence \answer{\text{the lines intersect at }(4,-1,-5)}.
\setcounter{enumi}{24}
\item \[-(x-5)+2(y+2)+3(z-4)=0,\qquad\answer{-x+2y+3z=3}.\]
\setcounter{enumi}{28}
\item Parallel planes have the same normal. Substitution of $(2.1,1.7,-0.9)$ gives \[\answer{2x-y+3z=-\frac15}.\]
\setcounter{enumi}{32}
\item Using $\langle1,-9,4\rangle\times\langle-4,-4,-1\rangle=5\langle5,-3,-8\rangle$, \[\answer{5x-3y-8z+9=0}.\]
\setcounter{enumi}{34}
\item A normal from the line direction and the vector to $(3,5,-1)$ is $\langle8,1,-2\rangle$, so \[\answer{8x+y-2z-31=0}.\]
\setcounter{enumi}{36}
\item A member of the pencil through the two given planes is $3P_1-4P_2=0$, yielding \[\answer{x-2y-z+3=0}.\]
\setcounter{enumi}{44}
\item Substitution into $x+2y-z=7$ gives $1+3t=7$, hence \answer{(-2,6,3)}.
\setcounter{enumi}{52}
\item The normal vectors have $|\vect n_1\cdot\vect n_2|=3$, so \[\answer{\theta=\cos^{-1}(1/\sqrt6)\approx65.9^\circ}\quad\text{(neither parallel nor perpendicular).}\]
\setcounter{enumi}{56}
\item Subtraction gives $y+z=0$ and $x=1$. Thus \[\answer{\vect r=\langle1,0,0\rangle+t\langle0,1,-1\rangle},\qquad \answer{\theta\approx15.8^\circ}.\]
\setcounter{enumi}{62}
\item The intercept form is \[\answer{\frac{x}{a}+\frac{y}{b}+\frac{z}{c}=1}.\]
\setcounter{enumi}{64}
\item A direction perpendicular to both $\langle1,1,1\rangle$ and $\langle1,-1,2\rangle$ is $\langle3,-1,-2\rangle$. Hence \[\answer{x=3t,\quad y=1-t,\quad z=2-2t}.\]
\setcounter{enumi}{70}
\item \[\answer{d=\frac{|3(1)+2(-2)+6(4)-5|}{\sqrt{3^2+2^2+6^2}}=\frac{18}{7}}.\]
\setcounter{enumi}{74}
\item Parallel planes with common normal $\langle a,b,c\rangle$ differ only by $d_1,d_2$. Projecting a point on one plane onto the unit normal gives \[\answer{D=\frac{|d_1-d_2|}{\sqrt{a^2+b^2+c^2}}}.\]
\end{enumerate}

\subsection*{12.6 Quadric Surfaces}
\begin{enumerate}[label=\textbf{\arabic*.}, start=3]
\item \[x^2+z^2=4\] is a circular cylinder of radius $2$ with \answer{\text{axis the }y\text{-axis}}.
\setcounter{enumi}{4}
\item $y=-x^2-1$ is a \answer{\text{parabolic cylinder}} with vertex line $(0,-1,z)$, opening in the negative $y$ direction.
\setcounter{enumi}{10}
\item For $x^2+y^2-z^2=1$, the $z=k$ traces are $x^2+y^2=1+k^2$ and the other coordinate traces are hyperbolas; it is a \answer{\text{hyperboloid of one sheet about the }z\text{-axis}}. Replacing it by $x^2-y^2+z^2=1$ rotates the axis to the $y$-axis. Completing the square in $x^2+y^2+2y-z^2=0$ gives $x^2+(y+1)^2-z^2=1$: the first surface translated one unit in negative $y$.
\setcounter{enumi}{12}
\item \[x=y^2+4z^2\] is an \answer{\text{elliptic paraboloid}} with vertex at the origin and axis in the positive $x$ direction.
\setcounter{enumi}{16}
\item \[\frac{y^2}{4}+\frac{z^2}{9}-\frac{x^2}{36}=1\] is a \answer{\text{hyperboloid of one sheet with axis the }x\text{-axis}}.
\setcounter{enumi}{20}
\item $y=z^2-x^2$ is a \answer{\text{hyperbolic paraboloid}} (saddle) with axis variable $y$.
\setcounter{enumi}{22}
\item The matches are identified as follows: $23$ is the ellipsoid with semiaxes $(1,1/2,1/3)$; $24$ is the ellipsoid with semiaxes $(1/3,1/2,1)$; $25$ is the one-sheet hyperboloid about $y$; $26$ is the two-sheet hyperboloid about $y$; $27$ is an elliptic paraboloid opening in $+y$; $28$ is an elliptic cone about $y$; $29$ is an elliptic cylinder parallel to $y$; and $30$ is the saddle $y=x^2-z^2$.
\setcounter{enumi}{32}
\item \[y^2=x^2+\frac{z^2}{9}\] is an \answer{\text{elliptic cone with axis the }y\text{-axis}}.
\setcounter{enumi}{44}
\item In cylindrical coordinates, the bounded region is \[\answer{1\le z\le2,\qquad1\le r\le z,\qquad0\le\theta\le2\pi};\] it lies between the cylinder $r=1$ and the cone $r=z$.
\end{enumerate}
